\subsection{RQ1: Effectiveness}
\label{sec:rq1-effectiveness}
\textbf{RQ1. How effectively does our design-centered approach generate complete project-level embedded software compared with existing LLM-based development approaches?}

\paragraph{Evaluation setup.}
For the RQ1 comparison, we select three projects from the nine collected real-world embedded software projects: Crazyflie, OnStep, and DiscoBot. Using C application tasks derived from these projects, we compare \textsc{EmbedDev} with MetaGPT, ChatDev, and StructGen under a shared API and three generation models: DeepSeek-V4-Pro, Gemini-3-Flash-Preview, and GPT-5.5. Three repeat identities per project give nine positions per method--model row: 108 shared-API positions, or 135 including the Direct-LLM reference. Direct generation without planning, repair, or the supplied API is included as a reference; design ablations appear in RQ2.

\begin{table*}[tbp]
\centering\footnotesize
\caption{RQ1 effectiveness by method and generation model. Entries are mean $\pm$ population standard deviation (percentage points; divisor 3) across three repeat blocks. Each block equally averages the three projects, retaining all nine project--repeat positions per row. Host requires compilation, linking, and all behavior scenarios. The two testcase columns use the static-assessment measures defined in the text. Requirement Coverage counts fully and partially covered requirement leaves. Bold scores denote the highest mean under the same model (ties included); Direct-LLM is a separate reference.}
\label{tab:rq1-method-model}
\setlength{\tabcolsep}{3pt}
\renewcommand{\arraystretch}{1.18}
\begin{tabular*}{\linewidth}{@{\extracolsep{\fill}}l|l|cccc@{}}
\toprule
\textbf{Method} & \textbf{Model} & \textbf{\shortstack{Host Success\\(\%) $\uparrow$}} & \textbf{\shortstack{Test Case\\Pass Rate (\%) $\uparrow$}} & \textbf{\shortstack{Test Case\\Pass Rate (H) (\%) $\uparrow$}} & \textbf{\shortstack{Requirement\\Coverage (\%) $\uparrow$}}\\
\midrule
Direct-LLM & DeepSeek-V4-Pro & $0.00 \pm 0.00$ & $0.00 \pm 0.00$ & $49.21 \pm 5.17$ & $81.42 \pm 14.68$\\
 & Gemini-3-Flash & $0.00 \pm 0.00$ & $0.00 \pm 0.00$ & $45.08 \pm 5.89$ & $66.11 \pm 15.20$\\
 & GPT-5.5 & $0.00 \pm 0.00$ & $0.00 \pm 0.00$ & $59.73 \pm 1.50$ & $87.88 \pm 1.49$\\
\midrule
MetaGPT* & DeepSeek-V4-Pro & $0.00 \pm 0.00$ & $0.00 \pm 0.00$ & $4.48 \pm 3.29$ & $13.89 \pm 9.83$\\
 & Gemini-3-Flash & $0.00 \pm 0.00$ & $0.00 \pm 0.00$ & $0.72 \pm 1.01$ & $2.78 \pm 3.93$\\
 & GPT-5.5 & $0.00 \pm 0.00$ & $0.00 \pm 0.00$ & $2.69 \pm 2.01$ & $15.53 \pm 11.01$\\
\midrule
ChatDev* & DeepSeek-V4-Pro & $88.89 \pm 15.71$ & $70.39 \pm 8.28$ & $75.50 \pm 1.26$ & $91.35 \pm 0.31$\\
 & Gemini-3-Flash & $44.44 \pm 15.71$ & $33.13 \pm 7.92$ & $61.99 \pm 6.75$ & $77.50 \pm 9.31$\\
 & GPT-5.5 & $77.78 \pm 15.71$ & $63.19 \pm 13.13$ & $68.19 \pm 13.61$ & $81.67 \pm 12.39$\\
\midrule
StructGen* & DeepSeek-V4-Pro & $33.33 \pm 0.00$ & $28.85 \pm 0.25$ & $63.48 \pm 1.91$ & $86.06 \pm 0.74$\\
 & Gemini-3-Flash & $22.22 \pm 15.71$ & $18.82 \pm 13.31$ & $55.13 \pm 8.94$ & $69.04 \pm 16.38$\\
 & GPT-5.5 & $55.56 \pm 15.71$ & $41.18 \pm 8.50$ & $70.18 \pm 1.49$ & $86.37 \pm 0.76$\\
\midrule
\textbf{EmbedDev} & DeepSeek-V4-Pro & $\mathbf{100.00 \pm 0.00}$ & $\mathbf{82.29 \pm 3.46}$ & $\mathbf{82.29 \pm 3.46}$ & $\mathbf{92.16 \pm 1.31}$\\
 & Gemini-3-Flash & $\mathbf{77.78 \pm 31.43}$ & $\mathbf{58.07 \pm 21.51}$ & $\mathbf{63.06 \pm 14.47}$ & $\mathbf{78.42 \pm 12.43}$\\
 & GPT-5.5 & $\mathbf{100.00 \pm 0.00}$ & $\mathbf{77.79 \pm 0.69}$ & $\mathbf{77.79 \pm 0.69}$ & $\mathbf{91.25 \pm 0.49}$\\
\bottomrule
\end{tabular*}
\par\vspace{3pt}
\begin{minipage}{\linewidth}\footnotesize
* Official implementations adapted to the C/API tasks. Direct-LLM receives no supplied API and serves as a separate reference. Gemini-3-Flash denotes Gemini-3-Flash-Preview. Both testcase columns are based on static assessment of unchanged code. The first is Host-gated; (H) additionally includes Host-failing artifacts through assisted code reading. These are not measurements of manually repaired code or executed testcase success. Requirement Coverage credits full and partial implementation equally. No-code positions contribute zero to the aggregate scores. Percentage-point differences use unrounded values where available. $\uparrow$ indicates higher is better.
\end{minipage}
\end{table*}
\paragraph{Measures.}
\textbf{Code delivery} records complete task-level code independently of internal acceptance. \textbf{Host success} requires compilation, linking, and all behavior scenarios: six for Crazyflie, six for OnStep, and eight for DiscoBot. These checks assess the application contract on a host computer, not upstream firmware on physical devices.

We additionally assess 100 Crazyflie, 90 OnStep, and 124 DiscoBot static testcases. For artifact $j$, its support is $q_j=P_j/(P_j+F_j+U_j)$, where $P$, $F$, and $U$ denote supported, failed, and uncertain items; inapplicable items are excluded. We report two complementary views:
\begin{equation}
S_{\mathrm{gated}}=\frac{100}{n}\sum_{j=1}^{n}H_jq_j,
\qquad
S_{\mathrm{static}}=\frac{100}{n}\sum_{j=1}^{n}\widetilde q_j.
\end{equation}
Here, $n=9$ is the number of project--repeat positions per method--model row, and $H_j=1$ only for Host-all-pass artifacts. In the static view, $\widetilde q_j=q_j$ for those artifacts; Host-failing code receives a fresh code-only LLM assessment without repair or added functionality. Explicit API and build-contract obligations remain assessable despite build obstacles. No-code positions contribute zero to both scores, retaining all nine positions in both denominators. Both measures concern static support; gating additionally requires executable integration.


\paragraph{Requirement coverage.}
We separately assess 342 requirement leaves derived from the frozen requirements of the three RQ1 projects: 145 for Crazyflie, 102 for OnStep, and 95 for DiscoBot. Of these, 318 are statically reviewable (141, 89, and 88, respectively). The remaining 24 require execution or hardware evidence, or clarification of the contract, and are excluded from this static metric. An LLM reads the unchanged candidate code and records a verdict, rationale, and source locations for each reviewable leaf: fully covered, partially covered, not covered, or uncertain. Partial coverage requires implementation evidence for part of the obligation; merely mentioning a requirement in a comment or function name is insufficient.

For candidate $j$, let $N_j$ be the fixed number of statically reviewable leaves, and let $F_j^{\mathrm{req}}$ and $P_j^{\mathrm{req}}$ count fully and partially covered leaves. Requirement coverage counts both categories:
\begin{equation}
c_j=\frac{F_j^{\mathrm{req}}+P_j^{\mathrm{req}}}{N_j}.
\end{equation}
Fully and partially covered leaves each receive one unit of credit; not-covered and uncertain leaves remain in the denominator without credit. The metric therefore measures the breadth of at least partial implementation. Host-failing code is assessed without repair, while no-code positions contribute zero only to the nine-position aggregates. We use the same equal-project repeat blocks and population standard deviation as the testcase measures. Requirement coverage is based on code-reading evidence and does not establish complete or dynamically verified requirement satisfaction.


\paragraph{Executable outcomes.}
Our approach delivers code at 9/9 DeepSeek, 8/9 Gemini, and 9/9 GPT positions. Of these 26 artifacts, 25 pass Host (Table~\ref{tab:rq1-method-model}). DeepSeek and GPT each pass 60/60 scenarios; Gemini passes 53/54 executed scenarios out of 60 planned. One OnStep position has no code, and one Crazyflie artifact passes five of six scenarios.

The Host-success rates are 100\%, 77.78\%, and 100\%, respectively, the highest observed under each model. ChatDev, the strongest Host baseline, passes 8/9, 4/9, and 7/9 positions, giving our approach margins of 11.11, 33.33, and 22.22 percentage points.

Code availability alone does not explain this result. StructGen delivers nine artifacts per model, but only 3/9 DeepSeek, 2/9 Gemini, and 5/9 GPT artifacts pass Host. Our delivery ties ChatDev and StructGen under DeepSeek, and StructGen under GPT; under Gemini it is lower than StructGen and ChatDev, despite higher Host success. MetaGPT delivers two DeepSeek, one Gemini, and two GPT artifacts, none Host-all-pass.

\paragraph{Testcase support beyond build success.}
Host-gated support is 82.29\% with DeepSeek, 58.07\% with Gemini, and 77.79\% with GPT, exceeding ChatDev, the strongest same-model baseline, by 11.91, 24.94, and 14.60 points. Both testcase scores coincide for our DeepSeek and GPT artifacts because all pass Host; Gemini's static score also credits the Crazyflie artifact failing one scenario.

Crediting Host-failing code raises StructGen's score from 28.85\% to 63.48\% with DeepSeek and from 41.18\% to 70.18\% with GPT: implemented functionality can coexist with unsuccessful integration. Our static scores remain highest, exceeding DeepSeek's strongest baseline, ChatDev (75.50\%), by 6.79 points and GPT's strongest baseline, StructGen (70.18\%), by 7.61 points.

\paragraph{Requirement-coverage results.}
The Requirement Coverage column of Table~\ref{tab:rq1-method-model} shows that \textsc{EmbedDev} has the highest aggregate requirement coverage under all three generation models: 92.16\% with DeepSeek, 78.42\% with Gemini, and 91.25\% with GPT. Its margins over the strongest other method or reference under each model are 0.80, 0.92, and 3.36 percentage points, respectively. These are descriptive aggregate differences across the three selected projects.

Direct-LLM also obtains relatively high requirement coverage (81.42\%, 66.11\%, and 87.88\%), although none of its artifacts passes every Host scenario. This illustrates that broad implementation of requirement content, including partial implementations, can coexist with integration failures. Requirement coverage complements the Host and testcase-support measures by assessing how many requirement obligations have implementation evidence.

\paragraph{Direct generation.}
Direct-LLM delivers 8/9 DeepSeek, 7/9 Gemini, and 9/9 GPT artifacts. None passes every Host scenario, although individual scenarios pass; static support is 49.21\%, 45.08\%, and 59.73\%. Partial functionality therefore coexists with integration failures. Because Direct-LLM omits the evaluated API, it is a separate reference, not a matched comparison of design modeling alone.

\paragraph{Interpretation.}
For each project--repeat position with code, we evaluate the last complete candidate from the selected execution. Rates describe the evaluated artifacts and do not estimate independent-attempt success or matched-budget efficiency. All static judgments have been manually reviewed without modifying the generated code. Static assessment complements execution-based validation.


% Preserve an existing rqanswer style; otherwise use a simple framed paragraph.
\ifcsname rqanswer\endcsname\else
\newenvironment{rqanswer}[1]{%
  \par\smallskip\noindent
  \begin{tabular}{|p{\dimexpr\linewidth-2\tabcolsep-2\arrayrulewidth\relax}|}
  \hline\textbf{Answer to RQ#1.}\ignorespaces
}{\\\hline\end{tabular}\par\smallskip}
\fi
\begin{rqanswer}{1}
Design-centered generation leads aggregate Host success, both testcase-support measures, and requirement coverage under all three models. Counting full and partial implementation, requirement coverage reaches 92.16\%, 78.42\%, and 91.25\% with DeepSeek, Gemini, and GPT, respectively. DeepSeek and GPT deliver and pass Host at all nine positions; Gemini delivers eight artifacts, seven Host-all-pass. These selected-artifact results support broader requirement coverage and stronger executable integration for the evaluated applications, while retaining residual delivery and behavioral failures.
\end{rqanswer}
